\chapter{Architecture}\label{ch:arch}

\section{Layers}

QuickEIP has three layers. Your application calls the public interface. The
portable core contains all protocol logic. The port connects the core to one
platform. Figure~\ref{fig:layers} shows the layers and the files in each.

\begin{figure}[htbp]
\centering
\begin{tikzpicture}[
  font=\sffamily\small,
  box/.style={draw=steel,fill=white,rounded corners=2pt,minimum height=0.62cm,inner sep=3pt,align=center},
  core/.style={box,fill=white,minimum width=2.95cm,text width=2.8cm,minimum height=0.8cm,font=\sffamily\scriptsize},
  port/.style={box,fill=teal!12,draw=teal,minimum width=2.25cm},
  shared/.style={box,fill=teal!28,draw=teal,minimum height=0.5cm,font=\sffamily\scriptsize},
  lbl/.style={font=\sffamily\scriptsize\bfseries\color{steel}}]

  % application
  \node[box,fill=navy,text=white,minimum width=12.6cm,minimum height=0.75cm] (app) at (0,6.75) {\bfseries Your application};

  % API headers
  \node[box,minimum width=6.1cm,text width=5.9cm,minimum height=0.95cm] (apiA) at (-3.25,5.55)
    {{\ttfamily\footnotesize eip.h}\\[-1pt]{\scriptsize scanner and built-in adapter}};
  \node[box,minimum width=6.1cm,text width=5.9cm,minimum height=0.95cm] (apiB) at ( 3.25,5.55)
    {{\ttfamily\footnotesize eip\_server.h\quad eip\_adapter.h}\\[-1pt]{\scriptsize your own adapter instance or event loop}};
  \draw[-Stealth,steel] (app.south -| apiA) -- (apiA);
  \draw[-Stealth,steel] (app.south -| apiB) -- (apiB);

  % core
  \node[lbl,anchor=west] at (-6.1,4.45) {SCANNER};
  \node[lbl,anchor=west] at ( 0.15,4.45) {ADAPTER};
  \node[font=\sffamily\scriptsize\color{steel},anchor=east] at (6.1,4.45) {\ttfamily src/\ \sffamily portable C99, no OS calls};
  \node[core] (s1) at (-4.72,3.75) {{\ttfamily\footnotesize eip\_session}\\[-1pt] tags, messages};
  \node[core] (s2) at (-4.72,2.80) {{\ttfamily\footnotesize eip\_io}\\[-1pt] Class 1 originator};
  \node[core] (s3) at (-1.57,3.75) {{\ttfamily\footnotesize eip\_scan}\\[-1pt] discovery};
  \node[core] (a1) at ( 1.57,3.75) {{\ttfamily\scriptsize eip\_adapter\_default}\\[-1pt] built-in instance};
  \node[core] (a2) at ( 1.57,2.80) {{\ttfamily\footnotesize eip\_server}\\[-1pt] sockets and timing};
  \node[core] (a3) at ( 4.72,2.80) {{\ttfamily\footnotesize eip\_adapter}\\[-1pt] sans-I/O core};
  \draw[-Stealth,steel!70] (a1) -- (a2);
  \draw[-Stealth,steel!70] (a2) -- (a3);
  \draw[-Stealth,steel!70] (s2) -- (s1);
  \node[core,minimum width=12.4cm,text width=12.2cm] (codec) at (0,1.85)
    {{\ttfamily\footnotesize eip\_codec}\quad packet encoding and parsing, used by both roles};
  \begin{scope}[on background layer]
    \node[draw=steel,dashed,rounded corners=4pt,fit=(s1)(s2)(s3)(a1)(a3)(codec),inner xsep=6pt,inner ysep=6pt,yshift=4pt,minimum height=3.6cm,fill=steel!5] (coreframe) {};
  \end{scope}

  % HAL bar
  \node[box,fill=steel,text=white,minimum width=12.6cm] (hal) at (0,0.65) {{\ttfamily eip\_hal.h}\quad port contract: 19 functions};
  \draw[-Stealth,steel] (codec.south) -- (hal.north);

  % ports
  \foreach \n/\x/\t in {pn/-5.2/netburner,pe/-2.6/esp-idf,pp/0/posix,pc/2.6/clearcore,pw/5.2/windows}{
    \node[port] (\n) at (\x,-0.55) {\ttfamily \t};
    \draw[steel] (hal.south -| \n) -- (\n.north);
  }
  \node[shared,minimum width=4.85cm] (bsd) at (-1.3,-1.45) {\ttfamily ports/bsd\_sockets};
  \node[shared,minimum width=2.25cm] (lw) at (2.6,-1.45) {\ttfamily ports/lwip\_raw};
  \draw[teal] (pe.south) -- (pe.south |- bsd.north);
  \draw[teal] (pp.south) -- (pp.south |- bsd.north);
  \draw[teal] (pc.south) -- (lw.north);
\end{tikzpicture}
\caption{The layers of QuickEIP. Arrows show the direction of calls. The
shared socket layers are reused by more than one platform port.}
\label{fig:layers}
\end{figure}

An application that has its own event loop can call the sans-I/O core in
\file{include/eip_adapter.h} directly. Most applications use \file{eip.h}
only.

\section{Source Files}\label{sec:sources}

Each build compiles three groups of files: the core, one port, and the
include paths. Table~\ref{tab:groups} lists them.

\begin{table}[htbp]
\centering
\caption{File groups in each build}\label{tab:groups}
\begin{tabularx}{\linewidth}{@{}P{0.16\linewidth}P{0.40\linewidth}L@{}}
\toprule
\textbf{Group} & \textbf{Files} & \textbf{Notes} \\
\midrule
Core & \file{src/eip_codec.c}, \file{src/eip_adapter.c},
\file{src/eip_session.c}, \file{src/eip_io.c}, \file{src/eip_scan.c},
\file{src/eip_server.c}, \file{src/eip_adapter_default.c}
& The same on all platforms. C99, no OS calls. \\
Port & One platform directory under \file{ports/}, plus a shared socket
layer for some platforms & Implements \file{include/eip_hal.h}. \\
Include paths & \file{include/} (public) and \file{src/} (internal headers
that the core needs) & Add both to the C and C++ compilers. \\
\bottomrule
\end{tabularx}
\end{table}

You can omit a core file if you do not use its functions.
Table~\ref{tab:corefiles} shows what each file is necessary for.

\begin{table}[htbp]
\centering
\caption{Core files and their functions}\label{tab:corefiles}
\begin{tabularx}{\linewidth}{@{}P{0.36\linewidth}L@{}}
\toprule
\textbf{File} & \textbf{Necessary for} \\
\midrule
\file{eip_codec.c} & All functions. \\
\file{eip_adapter.c}, \file{eip_server.c} & The adapter. \\
\file{eip_adapter_default.c} & \code{eip_adapter_start} and the other
built-in-instance calls. Holds a static \code{EipServer} of approximately
11~KB. \\
\file{eip_session.c} & Scanner sessions, tags, explicit messages, and
symbols. \\
\file{eip_io.c} & Scanner Class~1 I/O. Needs \file{eip_session.c}. \\
\file{eip_scan.c} & \code{eip_list_identity} and \code{eip_scan_devices}. \\
\bottomrule
\end{tabularx}
\end{table}

\begin{note}
\notelabel If you link with \code{-ffunction-sections -fdata-sections
-Wl,--gc-sections}, the linker removes the functions that you do not call.
Thus you can usually keep all the core files in the build.
\end{note}

\section{Execution Model}\label{sec:threads}

\subsection{The Adapter Context}

The adapter core runs in one context. On an RTOS port, this context is a task
that the port starts. On a bare-metal port, this context is your main loop,
which calls \code{eip_adapter_poll}. Table~\ref{tab:runmodes} compares the two
modes.

\begin{table}[htbp]
\centering
\caption{How the adapter runs}\label{tab:runmodes}
\begin{tabularx}{\linewidth}{@{}P{0.24\linewidth}P{0.26\linewidth}L@{}}
\toprule
\textbf{Mode} & \textbf{Ports} & \textbf{What you do} \\
\midrule
Task (RTOS) & NetBurner, ESP-IDF, POSIX, Windows & Nothing. \code{eip_adapter_start} starts a task. Do not call \code{eip_adapter_poll}. \\
Poll (bare metal) & ClearCore, other lwIP raw-API boards, Windows with \code{EIP_WIN_POLL_MODE} & Call \code{eip_adapter_poll} on each pass of the main loop, at intervals shorter than your fastest RPI. \\
\bottomrule
\end{tabularx}
\end{table}

\begin{caution}
Do not call \code{eip_adapter_poll} on an RTOS port. A task already runs the
adapter there. Two contexts that run the core at the same time can corrupt
the adapter state.
\end{caution}

\subsection{Rules for Tasks}

\begin{itemize}
  \item \code{eip_adapter_read}, \code{eip_adapter_write}, and
        \code{eip_adapter_snapshot} are safe from any task. They take
        \code{eip_hal_lock} for some microseconds. The lock is never held
        during a socket call.
  \item One task at a time can use a given \code{EipSession} or
        \code{EipIo}. Different sessions can run in different tasks.
  \item Scanner calls block until they are complete or until their timeout.
        On an RTOS, call them from a task of their own.
  \item On bare metal, a blocking scanner call keeps the network stack in
        operation, and the adapter sockets continue to receive packets. But
        the remainder of your main loop stops until the call returns.
  \item Do not call QuickEIP from an interrupt handler.
\end{itemize}

\section{Memory}\label{sec:memory}

QuickEIP uses no heap in the protocol code. Its RAM is static, or it is in
structures that you supply. Let $P$ be \code{EIP_MAX_PACKET} and let $C$ be
\code{EIP_ADAPTER_MAX_CLIENTS}. Then these approximations apply:
\begin{align}
  M_{\text{server}}  &\approx P\,(C + 3), \label{eq:server}\\
  M_{\text{session}} &\approx 3P. \label{eq:session}
\end{align}
With the defaults $P = 1536$ and $C = 4$, equation~\eqref{eq:server} gives
\num{10752} bytes, and equation~\eqref{eq:session} gives \num{4608} bytes plus
some fields. Table~\ref{tab:ram} gives the totals.

\begin{table}[htbp]
\centering
\caption{RAM with the default settings}\label{tab:ram}
\begin{tabularx}{\linewidth}{@{}P{0.38\linewidth}L@{}}
\toprule
\textbf{Item} & \textbf{RAM} \\
\midrule
Built-in adapter & Approximately 11~KB static, plus a 1.5~KB shared reply buffer. \\
Each \code{EipSession} & Approximately 4.7~KB. \\
Each \code{EipIo} & Approximately 100 bytes, plus a 1.5~KB stack buffer during produce and consume. \\
Scanner calls & Up to approximately 2~KB of caller stack. A tag write builds requests of 400 bytes. \\
Code, all functions & Approximately 24~KB at \code{-Os} on a Cortex-M4. \\
\bottomrule
\end{tabularx}
\end{table}

\begin{caution}
An \code{EipSession} is too large for most task stacks. Declare it
\code{static}, or put it in memory that you allocate one time at start-up.
A session on a small stack causes a stack overflow.
\end{caution}

The assembly data of the adapter comes from one pool. The sum of the sizes
$s_i$ of all assemblies must fit in it:
\begin{equation}
  \sum_{i=1}^{N} s_i \le \code{EIP_ADAPTER_DATA_POOL}, \qquad N \le \code{EIP_ADAPTER_MAX_ASSEMBLIES}.
  \label{eq:pool}
\end{equation}
